\input{../../../stock/en-tete_v6.tex}

\newcommand{\titrechapitre}{Rappels sur les asymptotes}
\newcommand{\numerochapitre}{1}

\renewcommand{\headrulewidth}{0.4pt}
\renewcommand{\footrulewidth}{0.4pt}
\fancyfoot[L]{Raphaël JONTEF}
\fancyfoot[C]{\thepage}
\fancyfoot[R]{\href{https://www.alveusclub.com/}{Alveus}}
\fancyhead[L]{Mathématiques : Terminale Générale}
\fancyhead[R]{\titrechapitre}

\begin{document}
\input{../../../stock/commands.tex}

\section*{Asymptotes d'une courbe en Terminale}

\subsection*{1. Asymptote horizontale}
Pour étudier une éventuelle asymptote horizontale en \(-\infty\) ou en \(+\infty\), on calcule :
\[
\lim_{x\to +\infty} f(x) \quad \text{ou} \quad \lim_{x\to -\infty} f(x)
\]
Si l'une de ces limites existent et vaut un réel \(\ell\), alors la courbe admet l'asymptote horizontale~:
\[
y=\ell.
\]
Une asymptote horizontale ne peut exister qu'en l'infini.

\subsection*{2. Asymptote verticale}
La droite \(x=a\) est une asymptote verticale lorsque :
\[
\lim_{x\to a^-} f(x)=\pm\infty
\quad\text{ou}\quad
\lim_{x\to a^+} f(x)=\pm\infty.
\]
Une asymptote verticale ne peut exister qu'en réel $a$.

\subsection*{3. Asymptote oblique (quasiment HP)}
Pour chercher une asymptote oblique en \(+\infty\) ou en \(-\infty\), on calcule :
\[
m = \lim_{x\to \pm\infty} \frac{f(x)}{x}.
\]
Si cette limite existe et \(m\neq 0\), alors on calcule :
\[
p = \lim_{x\to \pm\infty} \bigl(f(x)-mx\bigr).
\]
Si \(p\) est un réel, alors la courbe admet l’asymptote oblique :
\[
y = mx + p.
\]






\end{document}